\section{The variational theorem}\label{sec:var-theorem}

\subsection{The problem}\label{ss:25.1}

Let $V$ be a real configuration on $\T$. The \emph{variational problem} is to minimize the Wilson action $S$ on
\[
\mathfrak B(V)=\{U\in\mathfrak U_0(c_1)\ \text{real}:\ U^{(k)}=V\},
\]
the set of fine configurations of small curvature whose $k$-fold average is $V$ \cite[(3), (5), (6)]{BalVar}. All averages in this definition are defined: for $U\in\mathfrak U_0(c_1)$ the configurations $U^{(j)}$ lie in $\mathfrak U_j(2c_1)$, by Proposition~\ref{prop:avg-regular}, and for a real reference $U_0\in\mathfrak U_0(c_1)$ the relative configuration $UU_0^{-1}$ lies in the averaging domain up to scale $k$ (Section~\ref{sec:adapted}). So the domain of the group average is implied by the curvature bound $c_1$, through the condition $c_1\le r_M/C$ of Part~IV. Likewise, the small relative configurations $e^{\ii\eta A}$ of the Landau chart lie in the domain of the gauge-adapted averaging by Lemma~\ref{lem:one-step-expansion}, through $c_2,c_3\le r_M/C$. Both $S$ and $\mathfrak B(V)$ are invariant under the residual gauge group $\mathfrak G_0$. A \emph{critical orbit} is a $\mathfrak G_0$-orbit of critical points of $S$ restricted to $\mathfrak B(V)$. Its elements are Ba\l aban's \emph{background fields} (minimal configurations) $U_k(V)$.

Ba\l aban's Theorem~1 \cite{BalVar} solves this problem by induction over the number $k$ of averaging steps. The $k$-th step starts from a reference configuration $U_0$, the solution of the $(k-1)$-step problem for an extension of $V$. It solves the problem near $U_0$ by a contraction, and then re-establishes the regularity of the solution with constants independent of $k$, by a local analysis on cubes. The theorem below is the first part of this scheme: the problem relative to a given regular reference configuration \cite[Sects.~A--E, G]{BalVar}. The induction over $k$, with the local improvement of \cite[Sect.~F]{BalVar}, is carried out in Part~VI (Theorem~\ref{thm:balaban}).

\subsection{Statement}\label{ss:25.2}

Throughout this part, $U_0$ is a real reference configuration that is $(M,\eps)$-regular with $|J_{U_0}|_\infty\le\eps$. As before, $W_k=U_0^{(k)}$, and the operators of Parts~II and~III are taken at $U_0$ with the averaging pair $(Q'_k(U_0),\bQ_k(U_0))$. We fix once and for all a small constant $\theta_0>0$; the uniqueness statements refer to the class $\mathfrak U(\theta_0)$.

\begin{theorem}[The variational problem relative to a regular reference]\label{thm:variational}
There are $\theta_0,c_6,C_6,\gamma'>0$ such that the following holds. Let $U_0$ be as above, and let $V$ be a real configuration on $\T$ with
\[
V(c)W_k(c)^{-1}=e^{\ii B(c)},\qquad|B|_\infty\le\beta .
\]
If $\eps+\beta\le c_6$, then:
\begin{enumerate}[label=(\alph*)]
\item \emph{(Existence and uniqueness.)} $S|_{\mathfrak B(V)}$ has exactly one critical orbit in $\mathfrak B(V)\cap\mathfrak U(\theta_0)$.
\item \emph{(Minimality.)} This orbit is the set of minimizers of $S$ on $\mathfrak B(V)\cap\mathfrak U(\theta_0)$.
\item \emph{(Landau representative.)} The orbit contains $U_*=\mathsf U(A_*)=\big(e^{\ii\eta A_*}U_0\big)^{\mathsf w_k(e^{\ii\eta A_*})}$, where $A_*$ is real, $R(U_0)D^T_{U_0}A_*=0$, $\mathcal Q(e^{\ii\eta A_*})=e^{\ii B}$, and
\[
|A_*|_{C^1}+|\Lap_{U_0}A_*|_\infty+\sup_y\Hol{\nabla\tilde A_{*,y}}{B(y)}\le C_6(\eps+\beta).
\]
The proof gives the H\"older bound on the unit cubes $\Lam_1(c_y)$ in place of $B(y)$ (Theorem~\ref{thm:local-regularity}), and, by Remark~\ref{rem:cubes}, on every cube $\Lam_1(z)$. This is the form used in the proof of (e) and in Corollary~\ref{cor:regularity}.
To first order, $A_*=\bH B-\bC J_{U_0}+O\big((\eps+\beta)^2\big)$ in $|\cdot|_{C^1}$.
\item \emph{(Euler--Lagrange equation.)} There is a unique unit-lattice field $\omega$ with $\nabla\Phi(A_*)=d\mathcal Q(A_*)^T\omega$, where $\Phi(A)=S(e^{\ii\eta A}U_0)$ and $d\mathcal Q(A)$ is the derivative of $A\mapsto\ell(\mathcal Q(e^{\ii\eta A}))$. It satisfies $|\omega|_\infty\le C_6(\eps+\beta)$.
\item \emph{(Regularity of the minimizer.)} $U_*$ is $(M,\eps+C_6(\eps+\beta))$-regular and $U_*\in\mathfrak U(C_6(\eps+\beta))$.
\item \emph{(Analyticity and decay.)} $B\mapsto A_*(B)$ extends to an analytic map on complex unit-lattice fields with $\eps+|B|_\infty<c_6$, with $|A_*(B)|_{C^1}\le C_6(\eps+|B|_\infty)$. For complex $B$ the $C^1$ bound and the decay below are claimed, but not the second-order and H\"older bounds of (c). For real and complex $B$, if $\delta B$ is supported on the bond $c$, then
\[
\big|\,dA_*(B)\delta B\,(x)\big|+\big|\nabla\big(dA_*(B)\delta B\big)(x)\big|\le C_6e^{-\gamma'|x-c|}|\delta B| .
\]
\item \emph{(Covariance.)} For a real $v$ on $\T$, the critical orbit for $V^v$ is the image of the critical orbit for $V$ under any real $u$ with $u|_\T=v$.
\item \emph{(Lattice symmetry.)} Let $r\in\mathfrak S_k$. Then $rU_0$ satisfies the hypotheses above with the same $\eps$, $W_k(rU_0)=rW_k$, and $(rV)(c)\,(rW_k)(c)^{-1}=e^{\ii(r_{W_k}B)(c)}$ with $|r_{W_k}B|_\infty=|B|_\infty$. The critical orbit for $rV$ is the image under $r$ of the critical orbit for $V$. The Landau representatives and the multipliers satisfy
\[
A_*(r_{W_k}B;rU_0)=r_{U_0}A_*(B;U_0),\qquad\omega(r_{W_k}B;rU_0)=r_{W_k}\,\omega(B;U_0),
\]
where the second argument is the reference. The first identity also holds for complex $B$ as in~(f).
\end{enumerate}
\end{theorem}

Parts (a)--(c) and (e) are Ba\l aban's Theorem~1 \cite{BalVar}, relative to a reference and on a single torus; the regularity bounds are \cite[(9), (10)]{BalVar} with the reference-dependent constant. Parts (d) and (f) are \cite[(171), (178)]{BalVar} and \cite[Prop.~9, (190)]{BalVar}. Part~(h) is the lattice symmetry of the problem; here $r$ acts on configurations by $(rU)(b)=U(r^{-1}b)$, and $r_{U_0}$, $r_{W_k}$ act on Lie fields relative to $U_0$ and $W_k$ as in Proposition~\ref{prop:linearized-averages}(g). The constants do not depend on $k$, $n$ or $\eta$. They depend on $L$. The order in which they are fixed is given in the proof of the theorem, in Section~\ref{ss:27.6}. They may be taken independent of $M\ge M_0$, for $\eps$ below the threshold at $M=M_0$; this threshold includes $\eps_{\rm B}$ of Lemma~\ref{lem:uniformity}(c), taken at $M=M_0$. Lemma~\ref{lem:uniformity} is stated in Part~VII, but it, its proof and the two items before it (Definition~\ref{def:avg-parameter}, Lemma~\ref{lem:regular-monotone}) use only Parts~I--IV. Indeed, an $(M,\eps)$-regular reference is $(M_0,\eps)$-regular (Lemma~\ref{lem:regular-monotone}), the torus operators of Parts~II and~III do not depend on the regularity scale (Step~0 of the proof of Lemma~\ref{lem:uniformity}), and the proof of the regularity statement in (e) uses only local estimates on the cubes of Definition~\ref{def:regular}. The thresholds that depend on $M$ are those of Parts~VI--VIII, for instance $a_1$ of Theorem~\ref{thm:balaban} (whose constant $B_3$ also depends on $M$), $a_\delta$ of Proposition~\ref{prop:sharp-curvature} and $\eps_L\le\eps_{\rm cmp}(M'_L)$ with $M'_L\ge M$ (Theorem~\ref{thm:step-coercive}); the constants of Part~VII do not (Remark~\ref{rem:step-M}).

In (c), $A_*$ is the Landau representative of the orbit relative to $U_0$, and $U_*=\mathsf U(A_*)$ is its averaged axial representative. The configuration $e^{\ii\eta A_*}U_0$ is gauge equivalent to $U_*$, but it need not lie in $\mathfrak B(V)$, since $\mathsf w_k(e^{\ii\eta A_*})$ need not be $1$ on $\T$.

Part~(g) does not need (a): for real $u$ with $u|_\T=v$, the map $U\mapsto U^u$ is a bijection from $\mathfrak B(V)\cap\mathfrak U(\theta_0)$ onto $\mathfrak B(V^v)\cap\mathfrak U(\theta_0)$, because $\mathfrak a^{(k)}(U^u)=\mathfrak a^{(k)}(U)^{u|_\T}$. It preserves $S$ and critical points, and it maps $\mathfrak G_0$-orbits to $\mathfrak G_0$-orbits, since $u\mathfrak G_0u^{-1}=\mathfrak G_0$. Part~(h) is proved at the end of Section~\ref{sec:var-regularity}. The rest of this part proves (a)--(f). The proof has four steps:
\begin{enumerate}[label=\arabic*.]
\item Reduction to the Landau chart, by Corollary~\ref{cor:parametrization} (Section~\ref{subsec:reduction}).
\item Linearization of the constraint and solution by contraction (Section~\ref{sec:var-solution}).
\item Convexity, which gives minimality (Section~\ref{sec:var-solution}).
\item The Euler--Lagrange equation, which shows that the constructed point is critical on $\mathfrak B(V)$ and gives the second-order bounds (Section~\ref{sec:var-regularity}).
\end{enumerate}

\subsection{Reduction to the Landau chart}\label{subsec:reduction}

Let $\theta_0\le c_4/3$ and $\eps+\beta\le c_6\le c_4/3$, so that Theorem~\ref{thm:chart} and Corollary~\ref{cor:parametrization} apply with $\theta=\theta_0$. Let
\[
\mathcal M_V=\big\{A\ \text{real}:\ RD^TA=0,\ \ell(\mathcal Q(e^{\ii\eta A}))=B,\ |A|_{C^1}<\rho_1\big\},\qquad\rho_1=2B_1(\eps+\theta_0+\beta).
\]

\begin{lemma}\label{lem:reduction}
\begin{enumerate}[label=(\alph*)]
\item $\mathcal M_V$ is a smooth manifold, and $\Phi=S\circ\mathsf U$ on $\mathcal M_V$, with $\mathsf U$ as in Corollary~\ref{cor:parametrization}.
\item Every $U$ in a critical orbit in $\mathfrak B(V)\cap\mathfrak U(\theta_0)$ has an averaged axial representative $\mathsf U(A)$ with $A\in\mathcal M_V$, and $A$ is a critical point of $\Phi|_{\mathcal M_V}$.
\item For every $U\in\mathfrak B(V)\cap\mathfrak U(\theta_0)$ there is $A\in\mathcal M_V$ with $S(U)=\Phi(A)$.
\end{enumerate}
\end{lemma}

\begin{proof}
(a) The map $A\mapsto\ell(\mathcal Q(e^{\ii\eta A}))$, restricted to the slice $\{RD^TA=0\}$, has derivative $\bQ+d\Cav_k(A)$. This is onto: $\bQ\bH=I$, $\bH$ maps into the slice, and $|d\Cav_k(A)\bH|\le C\rho_1<1$ in the sup norm by Proposition~\ref{prop:Q-expansion}(b). Hence $\mathcal M_V$ is a submanifold. The identity $\Phi=S\circ\mathsf U$ holds by gauge invariance of $S$.

(b), (c) The orbit of $U$ has a unique averaged axial representative (Lemma~\ref{lem:axial}(c); note $U\in\mathfrak U(\theta_0)\subset\mathfrak U_0(c_1)$). By Corollary~\ref{cor:parametrization} it equals $\mathsf U(A)$ for some $A\in\mathcal S_V$, and by Theorem~\ref{thm:chart}, $|A|_{C^1}\le B_1(\eps+\theta_0+\beta)<\rho_1$, so $A\in\mathcal M_V$. For (b), let $A(t)$ be a smooth curve in $\mathcal M_V$ with $A(0)=A$. Since $\rho_1\le\frac43B_1c_4\le\frac23r_u$, the configurations $\mathsf U(A(t))$ are defined, real and smooth in $t$, and $\mathsf U(A(t))^{(k)}=\mathcal Q(e^{\ii\eta A(t)})W_k=V$. At $t=0$ the curvature is that of $U$, at most $\theta_0\le c_4/3<c_1$; by continuity, $\mathsf U(A(t))\in\mathfrak U_0(c_1)$ for small $|t|$. So $\mathsf U(A(t))$ is a curve in $\mathfrak B(V)$ through the critical point $\mathsf U(A)$, and $\frac{d}{dt}\Phi(A(t))\big|_{t=0}=\frac{d}{dt}S(\mathsf U(A(t)))\big|_{t=0}=0$.
\end{proof}

By Lemma~\ref{lem:reduction}(b), uniqueness in (a) follows once we show that $\Phi|_{\mathcal M_V}$ has at most one critical point. Existence requires the converse: the critical point of $\Phi|_{\mathcal M_V}$ that we construct gives a critical point of $S|_{\mathfrak B(V)}$. The converse is not a formal consequence of (b), because $\mathsf U$ involves the gauge transformation $\mathsf w_k$, which depends on $A$. It is proved in Section~\ref{sec:var-regularity} via the Euler--Lagrange equation.
